\documentclass[11pt]{article}
\usepackage[T1]{fontenc}
\usepackage[utf8]{inputenc}
\usepackage{lmodern}
\usepackage{amsmath,amssymb,amsthm,mathtools}
\usepackage[margin=1in]{geometry}
\usepackage{microtype}
\usepackage{xcolor}
\usepackage[colorlinks=true,linkcolor=blue!55!black,citecolor=blue!55!black,urlcolor=blue!55!black]{hyperref}

\newtheorem{theorem}{Theorem}
\newtheorem{proposition}{Proposition}
\newtheorem{remark}{Remark}
\newcommand{\R}{\mathbb R}
\newcommand{\BV}{\mathrm{BV}}
\newcommand{\TV}{J_0}
\newcommand{\W}{W_{0,\theta}}
\newcommand{\F}{F_\theta}
\newcommand{\ess}{\operatorname*{ess\,sup}}

\title{High-dimensional nonattainment for spatially weighted \(L^2\)--BV regularization}
\author{Solution packet for arXiv:1403.5579}
\date{11 August 2026}

\begin{document}
\maketitle

\begin{abstract}
Theorem 2.8 of Mazzieri--Spies--Temperini asserts existence and uniqueness of
a minimizer for a mixed spatially weighted \(L^2\)--BV functional in every
dimension under the assumptions \((1-\theta)^{-1}\in L^1\) and
\(\theta^{-1}\in L^\infty\).  We give an explicit counterexample in every
dimension \(n\ge3\).  A sequence of shrinking BV spikes fits the datum
exactly while both penalties tend to zero, so the infimum is zero and is not
attained.  We also prove the sharp dimensional repair: the stated hypotheses
do give existence and uniqueness for \(n\le2\).  The proof needs only the
continuous critical embedding \(\BV(\Omega)\hookrightarrow L^2(\Omega)\),
not the compact embedding claimed in the source.
\end{abstract}

\section{The source claim and the dimension gap}

Let \(\Omega\subset\R^n\) be bounded with Lipschitz boundary and let
\(0\le\theta\le1\).  The source defines
\[
 \W(u)=\sup_{\nu\in\mathcal V_\theta}
       \int_\Omega-u\,\operatorname{div}(\theta\nu)\,dx
\]
and studies, for \(\alpha_1,\alpha_2>0\),
\begin{equation}\label{eq:functional}
 \F(u)=\|Tu-v\|_Y^2
 +\alpha_1\int_\Omega(1-\theta)|u|^2\,dx
 +\alpha_2\W(u),\qquad u\in L^2(\Omega).
\end{equation}
Theorem 2.8 asserts that if
\begin{equation}\label{eq:source-hyp}
 \frac1{1-\theta}\in L^1(\Omega),
 \qquad \frac1\theta\in L^\infty(\Omega),
\end{equation}
then \eqref{eq:functional} has a unique global minimizer in
\(\BV(\Omega)\), with no restriction on \(n\) \cite{source}.  Its proof
obtains a BV bound and then invokes a ``compact embedding of
\(\BV(\Omega)\) in \(L^2(\Omega)\).''  Such an embedding cannot supply an
\(L^2\) bound when \(n\ge3\).

We use two elementary facts already available in the source:
\begin{equation}\label{eq:comparison}
 0\le \W(u)\le \TV(u),
 \qquad
 \TV(u)\le \|\theta^{-1}\|_\infty\W(u)
 \quad\text{if }\theta^{-1}\in L^\infty.
\end{equation}

\section{Counterexample in every dimension at least three}

\begin{theorem}\label{thm:counterexample}
For every integer \(n\ge3\) there are a ball \(\Omega\subset\R^n\), a
weight \(\theta\) satisfying \eqref{eq:source-hyp}, a bounded linear map
\(T:L^2(\Omega)\to\R\), and \(v\in\R\) such that \(\F\) has infimum zero
but no global minimizer.  Thus the existence assertion in source Theorem
2.8 is false in every dimension \(n\ge3\).
\end{theorem}

\begin{proof}
Fix \(0<R<1\), put \(\Omega=B(0,R)\), and set
\[
 \theta(x)=1-|x|^{n-1},
 \qquad
 \gamma=\frac{n+2}{4},
 \qquad
 g(x)=|x|^{-\gamma}.
\]
Because \(n\ge3\),
\[
 1<\gamma<\frac n2.
\]
Consequently \(g\in L^2(\Omega)\), and
\[
 Tu:=\int_\Omega g(x)u(x)\,dx
\]
defines a bounded linear functional on \(L^2(\Omega)\).  Take \(v=1\).
The weight has
\[
 \theta(x)\ge1-R^{n-1}>0,
 \qquad
 \frac1{1-\theta(x)}=|x|^{-(n-1)}\in L^1(\Omega),
\]
so both source hypotheses hold.

Let \(\sigma_{n-1}=|S^{n-1}|\).  For \(0<r<R\), define
\[
 c_r=\frac{n-\gamma}{\sigma_{n-1}}r^{\gamma-n},
 \qquad
 u_r=c_r\mathbf1_{B(0,r)}.
\]
Each \(u_r\) belongs to \(L^2(\Omega)\cap\BV(\Omega)\), and radial
integration gives
\begin{equation}\label{eq:datafit}
 Tu_r=c_r\sigma_{n-1}\int_0^r\rho^{n-1-\gamma}\,d\rho=1.
\end{equation}
Thus the data term in \eqref{eq:functional} vanishes identically along this
sequence.  The weighted quadratic term is
\begin{align}
 \int_\Omega(1-\theta)|u_r|^2\,dx
 &=c_r^2\sigma_{n-1}\int_0^r\rho^{2n-2}\,d\rho\notag\\
 &=\frac{(n-\gamma)^2}{\sigma_{n-1}(2n-1)}
   r^{2\gamma-1}
 =\frac{(n-\gamma)^2}{\sigma_{n-1}(2n-1)}r^{n/2}.
 \label{eq:weighted}
\end{align}
Moreover, by \eqref{eq:comparison} and the perimeter formula for a ball,
\begin{equation}\label{eq:variation}
 \W(u_r)\le\TV(u_r)
 =c_r\sigma_{n-1}r^{n-1}
 =(n-\gamma)r^{\gamma-1}
 =(n-\gamma)r^{(n-2)/4}.
\end{equation}
Both exponents in \eqref{eq:weighted}--\eqref{eq:variation} are positive.
Therefore
\[
 0\le\F(u_r)
 \le \alpha_1\frac{(n-\gamma)^2}{\sigma_{n-1}(2n-1)}r^{n/2}
 +\alpha_2(n-\gamma)r^{(n-2)/4}
 \longrightarrow0.
\]
Hence \(\inf_{L^2}\F=0\).

If a minimizer \(u_*\) existed, all three nonnegative terms in
\eqref{eq:functional} would have to vanish.  Since
\(1-\theta=|x|^{n-1}>0\) almost everywhere, vanishing of the weighted
quadratic term forces \(u_*=0\) almost everywhere.  But then
\(Tu_*=0\ne1=v\), contradicting vanishing of the data term.  Thus the
infimum is not attained.
\end{proof}

\begin{remark}
Taking the zero set \(\{\theta=0\}\) to be empty shows that the same example
also contradicts the unrestricted-dimensional existence assertion in source
Theorem 2.10, whose hypotheses then specialize to those of Theorem 2.8.
\end{remark}

\section{The sharp repair in dimensions one and two}

\begin{proposition}\label{prop:repair}
Under the source hypotheses \eqref{eq:source-hyp}, the conclusion of source
Theorem 2.8 is valid when \(n\le2\): the functional \(\F\) has a unique
global minimizer in \(\BV(\Omega)\).
\end{proposition}

\begin{proof}
Consider a sequence \((u_k)\) on which the two penalty terms are bounded.
Cauchy--Schwarz and \((1-\theta)^{-1}\in L^1\) give
\[
 \|u_k\|_{L^1}
 \le \left\|\frac1{1-\theta}\right\|_{L^1}^{1/2}
     \|\sqrt{1-\theta}\,u_k\|_{L^2}.
\]
The second inequality in \eqref{eq:comparison} gives a uniform bound on
\(\TV(u_k)\).  Hence \((u_k)\) is bounded in \(\BV(\Omega)\).  For
\(n=2\), the critical Sobolev inequality gives the continuous embedding
\(\BV(\Omega)\hookrightarrow L^2(\Omega)\); for \(n=1\), BV embeds into
\(L^\infty\), hence also into \(L^2\).  Thus \((u_k)\) is bounded in the
reflexive space \(L^2(\Omega)\) and has a weakly convergent subsequence.

The source proves weak lower semicontinuity of \(\W\); the other terms are
weakly lower semicontinuous as well.  The direct method therefore produces
a minimizer.  The reciprocal assumption implies \(1-\theta>0\) almost
everywhere, so \(u\mapsto\|\sqrt{1-\theta}\,u\|_2^2\) is strictly convex.
This gives uniqueness, while the displayed \(L^1\) and variation bounds put
the minimizer in \(\BV(\Omega)\).
\end{proof}

\section*{Proof Intuition}

The condition \((1-\theta)^{-1}\in L^1\) controls the \(L^1\) mass of a
candidate, and the weighted variation controls its BV seminorm.  In
dimensions one and two, those two controls automatically bound the
\(L^2\) norm, so weak \(L^2\) compactness completes the direct method.  In
dimension three and higher, BV is subcritical for \(L^2\): a shrinking ball
can carry taller and taller spikes while its perimeter cost vanishes.  By
placing the degenerating quadratic weight at the same point and choosing a
singular but square-integrable observation functional, the spikes fit the
datum exactly.  The objective approaches zero, although zero cost itself
would force the candidate to be both \(0\) and data-fitting.  This is the
precise compactness failure hidden by the dimension-free statement.

\begin{thebibliography}{9}
\bibitem{source}
G. L. Mazzieri, R. D. Spies, and K. G. Temperini,
\emph{Mixed spatially varying \(L^2\)--BV regularization of inverse
ill-posed problems}, arXiv:1403.5579 (2014), especially Theorem 2.8 and its
proof on pp. 7--8.
\end{thebibliography}

\end{document}
